\section{Analytic certificates: leaving the polynomials}
\label{sec:analytic}

Everything up to here has been polynomial or rational. Cones, quadratics,
Bernstein trees, Sturm chains, closures, and free-algebra words all manipulate
exact rational data, and the reason is the reason given in
Section~\ref{sec:trust}: exact data is what a checker can re-derive without
tolerance. An expression containing $e^x$ is not that.

Four contributing efforts take the step anyway, and they converge on the same
answer to the representation question. The input class is the \emph{rational
exponential polynomial}
\begin{equation}\label{eq:exppoly}
 f(x)=\sum_{\lambda\in\Lambda}P_\lambda(x)\,e^{\lambda x},
 \qquad \Lambda\subset\Q\ \text{finite},\quad P_\lambda\in\Q[x],
\end{equation}
which is exact rational data --- a finite list of rates and coefficient vectors
--- denoting a transcendental function. The goal is $f(x)\ge0$ on a ray or an
interval. Nothing here is approximated.

\subsection{The one lemma all of them prove}

The mechanism is an integrating factor. Multiplying by $e^{-cx}$ and
differentiating kills one exponential, and the key point is that this operation
stays inside the class \eqref{eq:exppoly}:

\begin{lemma}[Integrating-factor step]\label{lem:factorstep}
For $f$ of the form \eqref{eq:exppoly} and $c\in\Q$, set
$T_cf=\bigl(e^{-cx}f\bigr)'e^{cx}=f'-cf$. Then $T_cf$ is again of the form
\eqref{eq:exppoly}, with rate set contained in $\Lambda$, and if $f(a)\ge0$
and $T_cf\ge0$ on $[a,\infty)$ then $f\ge0$ on $[a,\infty)$.
\end{lemma}
\begin{proof}
The class is closed under differentiation and under scalar multiples, so
$f'-cf$ lies in it and introduces no new rate. For the inequality, let
$g(x)=e^{-cx}f(x)$. Then $g'(x)=e^{-cx}(T_cf)(x)\ge0$ on $[a,\infty)$, so $g$
is nondecreasing there and $g(x)\ge g(a)=e^{-ca}f(a)\ge0$. Since $e^{cx}>0$,
$f(x)=e^{cx}g(x)\ge0$.
\end{proof}

Three of the four efforts prove this lemma independently. Two of them state it
in literally the same form up to the names of the variables; the third proves
it at an arbitrary rational anchor $a$, of which the other two are the case
$a=0$. It is stated once here. Iterating it gives a \emph{ladder}: a sequence of
cofactors $c_1,\ldots,c_k$ reducing $f$ step by step, with one initial-value
check per rung.

That much is common ground. What the efforts do \emph{not} share is where the
ladder stops, and that turns out to be the whole design question.

\subsection{Three ladders, three terminal conditions}

\paragraph{Terminate at zero, and allow new rates.} \src{r4} drives the ladder
until the function is identically zero, and permits cofactors \emph{outside} the
rate set of the input --- auxiliary rates below the spectrum, which it calls
ghosts. Adding a rate enlarges the reachable language, so this accepts strictly
more inputs than a ladder confined to $\Lambda$; the price is that the ghost
search is a search, and its recorded ablation shows the factor count climbing as
the ghost parameter shrinks. It also proves a dominance theorem: cofactors in
ascending order succeed whenever any order does, which turns an apparent
$k!$ choice into one canonical attempt.

\paragraph{Terminate at a nonnegative polynomial, and fix the alphabet.}
\src{r6} restricts cofactors to exactly the input rates with exactly the
annihilator multiplicities, and stops as soon as the remaining function is an
ordinary polynomial with nonnegative coefficients --- which is visibly
nonnegative on $x\ge0$. Because the alphabet is fixed, it can go further than
\src{r4} on the structural side: it proves an \emph{if and only if} feasibility
criterion, a forced-multiplicity lemma, a monotone zero-count lemma, and a
shortest-ladder compiler that emits a separately checkable minimality receipt.

These two are not competing proofs of one theorem, and the relationship is worth
stating exactly because it is easy to get backwards. \src{r6}'s own completeness
lemma says any of its ladders extends to one whose terminal is the zero
function. So on the common domain --- anchor $0$, cofactors drawn from the input
spectrum --- the two accept the same functions, and \src{r6}'s early stopping is
a \emph{certificate-length} optimisation rather than a strength increase. The
strength increase is on \src{r4}'s side, and it comes from the ghosts, not from
the terminal condition.

\begin{remark}[Same argument, twice, for different purposes]
The two efforts also both prove that adjacent cofactors may be exchanged, by
the same algebra: $T_dT_ch=T_cT_dh+(c-d)\,T_ch$ evaluated at the anchor. One
uses it to justify sorting the cofactors ascending; the other to drive a
minimality argument. Neither cites the other. The identity is stated once in
this document and used for both.
\end{remark}

\paragraph{Do not fix a terminal at all.} \src{r5} runs an unstructured
breadth-first search over cofactors with a wider terminal cone and caps on depth
and node count. It has no ordering theorem, no feasibility criterion and no
minimality result, and it is the weakest of the three ladders --- a fact visible
in its own corpus, where the ladder family is a small minority of the records.
Its contribution is elsewhere, and is genuinely its own: see below.

\subsection{What the ladder cannot reach}

Two boundaries are worth naming, because both are sharp and both were found by
the efforts themselves rather than by a reviewer.

The first is the \emph{rate field}. Every rate in \eqref{eq:exppoly} must be
rational. The clearest statement of the consequence is that $2^x$ has rate
$\log 2$ and is therefore \emph{outside the fragment} --- not approximately
inside it, outside it. A rational approximation to $\log 2$ does not give an
exact source translation, and one of the efforts says so in exactly those
terms. The same applies at the anchor: at a general rational $a$ the values
$e^{\lambda a}$ need not be rational, so an initial-value check performed on
coefficients alone would be checking the wrong thing.

The second is the \emph{zero}. A ladder certifies $f\ge0$; it does not handle a
target that touches zero in the interior, and neither does any interval
enclosure, because an enclosure that straddles zero is inconclusive there no
matter how tight it is. Increasing precision cannot fix this --- it is not a
precision problem. The answer in the collection is exact vanishing information:
a \emph{jet} certificate recording the order of vanishing and the sign of the
first nonzero derivative, which settles the sign near the zero by an argument
that never divides by a quantity that vanishes. \src{r3} supplies this lane and
pairs it with a bounded automatic anchor/order selector.

\subsection{Remainders that keep the order of vanishing}

\src{r5}'s own contribution is a different object: anchored rational remainder
certificates for elementary functions --- $\exp$, but also $\sin$, $\cos$,
$\log(1+ax)$ and $\arctan$ --- which bound the tail of a truncated series while
\emph{preserving the order at which the function vanishes}. That last property
is what makes them composable with the jet lane instead of destroying the
information it needs.

The implementation detail that matters here is a Lean detail, and the effort
states it as a design rule rather than discovering it later: the certificate
never divides. A factorisation $f(x)=x^3g(x)$ is an identity that covers $x=0$
directly, whereas rewriting $g(x)=f(x)/x^3$ would need a nonzero side condition
that is easy to carry incorrectly into a total-division setting. The trap is
avoided structurally rather than by remembering to discharge a guard.

Alongside the remainders, the same effort supplies convergence and divergence
certificates for rational \emph{series}: an upper barrier with a Cauchy-modulus
witness on one side, and a harmonic-comparison divergence proof with a
least-index witness on the other.

\subsection{What the checkers actually verify, and one thing none of them says}

Every certificate in this section is a list of rationals, and every checker
re-derives the identity it claims. A ladder receipt is checked by recomputing
each $T_cf$ from the \emph{supplied} $f$ and comparing coefficient vectors, and
by evaluating the initial-value condition at the anchor exactly; a minimality
receipt is checked separately from the ladder it annotates, so a false
minimality claim cannot make an unsound ladder look sound. An enclosure is
checked by expanding it, not by sampling it. Where the efforts run
floating-point or high-precision numerical comparisons at all, they mark the
output files as diagnostics and keep them off the acceptance path --- one such
file is literally named to say so.

There is, however, a gap common to all four, and it is the same gap
Section~\ref{sec:closure}'s summation worker hit from the other side. Lean's
division is total: $x/0=0$. None of these four efforts names that convention.
They avoid its consequences --- by attaching positivity side conditions to
denominators, by factoring instead of dividing, by demanding \code{HasDerivAt}
witnesses rather than relying on a totalised \code{deriv} --- and the avoidance
is correct in each case. But a reader assembling these workers into a Lean
implementation should know that the discipline is being followed rather than
stated, and that Section~\ref{sec:closure}'s remark on total division applies
here with more force, not less.
